Pemetaan
\maabbdisp {u} { I } {\R^n
} {}
dinyatakan terhadap penampang basis
\mathl{e_1 , \ldots , e_n}{}
oleh
\mavergleichskettedisp
{\vergleichskette
{ u
}
{ =} { u_1e_1 + \cdots + u_ne_n
}
{ } {
}
{ } {
}
{ } {
}
}
{}{}{}
.
Menurut
\definitionsverweis {turunan vertikal}{}{}
dari koneksi tersebut, turunan penampang ini dalam satu-satunya arah turunan $\partial$, berdasarkan
Lema,
adalah

\mavergleichskettedisp
{\vergleichskette
{  \sum_{k = 1}^n \left(    \partial u_k  +  \sum_{j = 1}^n u_j \Gamma_{j}^k  \right) e_k
}
{ =} { \sum_{k = 1}^n \left(   u_k'  +  \sum_{j = 1}^n u_j \Gamma_{j}^k  \right) e_k
}
{ } {
}
{ } {
}
{ } {
}
}
{}{}{.}
Menurut
Lema,
$u$ merupakan penampang horizontal jika dan hanya jika turunan vertikalnya sama dengan $0$. Untuk setiap komponen, ini berarti

\mavergleichskettedisp
{\vergleichskette
{ u_k' + \sum_{j = 1}^n \Gamma_j^k u_j
}
{ =} { 0
}
{ } {
}
{ } {
}
{ } {
}
}
{}{}{}
untuk semua $k$. Kondisi ini ekuivalen dengan

\mavergleichskettedisp
{\vergleichskette
{ u_k'
}
{ =} { \sum_{j = 1}^n a_{kj} u_j
}
{ } {
}
{ } {
}
{ } {
}
}
{}{}{}
untuk semua $k$, yang tepat berarti bahwa $u$ memenuhi sistem matriks

\mavergleichskettedisp
{\vergleichskette
{ u'
}
{ =} { \left(  a_{kj}    \right) u
}
{ } {
}
{ } {
}
{ } {
}
}
{}{}{}
;
dengan demikian, pemetaan itu merupakan suatu solusi sistem persamaan diferensial linear tersebut.
